\section{Computations}\label{app:comp}

This appendix describes the finite computations. Every result marked as computer-assisted in the text is the output of a terminating exact search over a finite set; where random sampling is used, this is said, and such runs serve only as checks. The directory \texttt{code/} contains the programs, with instructions for running them and their expected outputs, and the logs of the long runs; the program files are named below next to each computation. Solution tables and other data files of more than $5$\,MB are not included; the programs that produce them are.

\subsection{The censuses}\label{app:census}

The censuses of Theorems~\ref{thm:110} and~\ref{cl:thm:census}, and the control rows of Table~\ref{tab:census} for $54$R, $90$R and $150$R, use the encoding, the enumeration, the re-verification and the invertibility test of Section~\ref{sec:method}. The programs and their logs are in \texttt{code/census}; the raw solution tables are not included, and \texttt{enum\_r4.py} produces those of radius $4$ again.

\emph{Coverage of Theorem~\ref{cl:thm:census}.} The enumeration (\texttt{code/census/centralizer.py}, with a cap of $5000$ solutions) is run at $R=1,2,3$ on $72$ rules, which represent $55$ of the $57$ classes without a fixed word of length at most $3$: all except the classes of $223$ and $255$, whose status is decided otherwise. Each class is represented by its smallest rule, and the classes of $13$, $14$, $15$, $25$, $26$, $28$, $29$, $30$, $41$, $43$, $57$, $58$, $62$, $75$, $77$, $94$ and $105$ also by the rules $162$, $42$, $170$, $152$, $74$, $56$, $184$, $106$, $134$, $142$, $156$, $78$, $110$, $154$, $178$, $122$ and $150$; this suffices by Lemma~\ref{lem:classes}. The summary lines of all runs are in \texttt{code/census/logs/census\_R123\_summary.txt}. For all $28$ classes of status (O), and for the classes of $14$, $46$, $63$, $79$, $94$, $95$, $110$, $133$ and $187$, the enumeration finishes below the cap at every $R\le3$, so the enumeration is complete there. For $133$ the solutions at $R\le3$ are the maps $\sh^aF^b$ only. For the classes of $30$, $45$, $75$ and $135$ the solutions at $R\le3$ are the maps $\sh^aF^b$ and the constants; for $24$ and $189$ there are in addition $5$, $13$, $25$ invertible maps at $R=1,2,3$, the maps $\kappa\sh^aF^b$; for $63$ and $187$ there are in addition $4$, $12$, $24$ invertible maps, the radius-deficient powers; for $27$ the solutions at $R\le2$ are the $13$ maps $\sh^aF^b$ and one constant map. The cap is reached only for classes whose status is decided otherwise: those of $59$, $127$, $179$, $90$, $15$ (run as $15$ and $170$), $60$, $153$, $165$, $126$ and $219$. For $27$ ($1742$ solutions at $R=3$) and for $105$ and $150$ ($4096$ solutions each at $R=3$) the enumeration finishes below the cap, and the solutions at $R=3$ are not classified; these classes have status (G) and (A), and Theorem~\ref{cl:thm:census} does not concern them. Tables~\ref{cl:tab:finite}--\ref{cl:tab:open} are assembled from the conjugacy classes, the fixed words of length at most $3$, the tests for affine and outer-permutive rules and the search over the $23$ nontrivial cell permutations (\texttt{code/census/classes88.py}), the hidden-bit and time-periodic walls of \cite{PaperC}, and the census output.

\emph{Radius $4$.} The census at radius $4$ of Section~\ref{cl:sec:holdouts} uses the same encoding with MiniSat~2.2 and a cap of $1000$ solutions (\texttt{enum\_r4.py}; summary lines in \texttt{logs/census\_R4\_summary.txt}), on one rule of each of the $28$ open classes and on the rules $14$, $24$, $46$, $62$, $63$, $79$, $94$, $126$, $133$, $187$, $189$ and $219$. Every solution is rechecked on all $4^{11}$ windows and matched against the $41$ maps $\sh^aF^b$ and the constants; injectivity is not tested at this radius. For the open classes and for $14$, $46$, $62$, $79$ and $94$ there is no further solution; for $133$ there are $27$, and for $126$ and $219$ the cap is reached. The $41$ further solutions for $24$ and $189$ are the maps $\kappa\sh^aF^b$, and the $40$ for $63$ and $187$ are compared exactly with the tables of the powers $\sh^aF^b$ whose neighbourhood lies in $[-4,4]$ (\texttt{identify\_r4.py}). The log also contains a run for rule $95$, which reproduces the count of $262$ solutions of \cite[Lemma~6.20]{PaperE}.

\emph{Cross-check.} The program \texttt{crosscheck.py} repeats the enumerations with other solvers, compares the solution sets with one another as sets of tables and with the logs of the enumerations (numbers of clauses and of solutions, the counts of the maps $\sh^aF^b$, of the constants and of the others, and the set itself whenever the log determines it), and rechecks every table on all windows. At $R\le3$ it covers the $72$, $67$ and $61$ instances at $R=1,2,3$ that finish below the cap, with Glucose~4, MiniSat~2.2 and CaDiCaL~1.5.3 (log \texttt{logs/crosscheck\_R123.txt}); at $R=4$ the rules $110$ and $30$ with Glucose~4 and CaDiCaL~1.5.3 (\texttt{logs/crosscheck\_R4\_110\_30.txt}) and the other $39$ rules that finish below the cap with Glucose~4 (\texttt{logs/crosscheck\_R4.txt}). There is no disagreement. With the option \texttt{--rup} the formula consisting of the encoding and one blocking clause for each solution is solved by Glucose~4 with proof logging, and the proof is checked lemma by lemma by reverse unit propagation, by a checker of about $70$ lines in the same file; this was done for all instances at $R\le3$ except those of the rules $27$, $105$ and $150$ at $R=3$, which have $1742$, $4096$ and $4096$ solutions. The cross-check was run on a machine with 94\,GB of memory, on which CaDiCaL solves the instances of radius $4$.

\emph{Rule $90$ at radius $2$.} For the control row of Table~\ref{tab:census}, the invertibility of each of the $2^{17}$ commuting maps was tested with the pair-graph test, in $5{,}851$ seconds: $26$ are invertible, all affine (\texttt{linear\_90.py}, \texttt{analyze\_90.py}; log \texttt{logs/analyze\_90\_R2.txt}).

\subsection{Computer-assisted statements}\label{app:stated}

\emph{Walls and searches} (Section~\ref{cl:sec:holdouts} and Proposition~\ref{cl:prop:nowalls}; programs and logs in \texttt{code/searches}). The searches are run on one rule of each of the $28$ open classes and on the rules $14$, $46$, $62$, $63$, $79$, $94$, $95$, $110$, $133$ and $187$. \emph{Fixed and time-periodic words, families of type (W1)} (\texttt{w1.c}): for $(K,p)=(11,1),(9,2),(7,3),(5,4),(3,5)$ and all $|s|\le p$, the graph on $\A^K$ with an edge from each word to each of its images under $\sh^sF^p$ in all contexts is built, and its strongly connected components that no edge leaves are computed; a proper invariant family, in particular a word fixed by $\sh^sF^p$ in every context, exists if and only if there is such a component other than $\A^K$. There is none, except at $s=0$ for $46$ and $191$ (level sets of their frozen fields) and for $133$ (its wall). For $p=1$, $s=0$ and $K\le6$ this is contained in the examination of local conservation laws in \cite[Section~4]{Takesue1989}; the frozen field of $46$ is the local conservation law of \cite[(4.4)--(4.6)]{Takesue1989}, that of $191$ is recorded in \cite[Table~3]{HT1991}, and the wall of $133$ is that of \cite[p.~391]{Takesue1989}. \emph{Hidden-bit walls} (\texttt{hb.c}): for $K=10$ and every position $1\le j\le8$, the largest set of pairs of words that differ exactly in the cell $j$ and whose images under $F$ in every context are again such pairs is computed; it is empty except for $133$. \emph{Eigenfields} (\texttt{factor.c}): for $(r,p)=(5,1),(4,2),(3,3),(2,4),(1,5)$ and all $|k|\le p$, the finest partition of $\A^{2r+1}$ compatible with $\vartheta\circ F^p=\sh^{-k}\circ\vartheta$ is computed by union--find, and likewise for binary $\vartheta$ with $\vartheta\circ F^p=\neg\,\sh^{-k}\circ\vartheta$; a nonconstant solution exists only at $k=0$, for $46$, $133$ and $191$. \emph{Conserved quantities} (\texttt{conserved.c}, \texttt{cq.py}): the additive conserved quantities of $F_f$ with a density $h$ on $r$ adjacent cells are the solutions of the linear system $h(F(w))-h(w_1\cdots w_r)=g(w_1\cdots w_{r+1})-g(w_0\cdots w_r)$, $w\in\A^{r+2}$, in the unknowns $h$ and $g$, where $F(w)$ is the image of the $r$ inner cells; it is solved for densities on $r\le5$ adjacent cells, modulo the prime $32749$ (which bounds the dimension over $\mathbb Q$ from above) and for $r\le3$ over $\mathbb Q$; nontrivial densities exist for $26$, $46$, $77$, $91$, $94$, $95$, $123$ (from $r=2$), $191$ (from $r=3$) and $133$, $175$ (from $r=4$), and for no other rule of the list. For densities on at most three adjacent cells this agrees with \cite[Table~3]{HT1991} and \cite[Table~5]{Takesue1995}, which do not go further; it is consistent with \cite{Takesue1997}. For $110$R (\texttt{nowalls110.py}), the right end cell of a word is fixed by $F^P$ in every context for no word and no $P\le6$.

\emph{Reach of the edge-move method} (Proposition~\ref{cl:prop:reach}, Remark~\ref{cl:rem:reach}).\label{app:reach} For the proposition, for each of the $44$ rules of the $11$ classes, all spatially periodic configurations of spatial period at most $6$ are enumerated with their $F$-orbits, and those are kept that satisfy condition (1) of Lemma~\ref{cl:lem:stuckcols} at some pair of adjacent columns in every phase of the orbit; this includes every background with a stuck right edge. For $1\le b\le4$ and $0<|a|\le b$, the window graph of $\mathrm{Fix}(\sh^aF^b)$ has as vertices the words of length $2b$ and as edges the words $w$ of length $2b+1$ whose image under $F^b$, a single cell, equals the cell $w_{b+a}$; it is pruned to the vertices on bi-infinite paths. For every kept background $B$ in $\mathrm{Fix}(\sh^aF^b)$, the program checks that every window of $B$ has exactly one outgoing edge (\texttt{code/reach/reach1\_unique.py}); a second program computes the strongly connected components and checks that no path from a window of $B$ reaches a recurrent component other than that of $B$ (\texttt{reach1.py}, which needs SciPy). Both find the known kinks for the classes of $110$, $14$ and $79$, and kinks for $25$, $31$ and $159$. For the remark (\texttt{code/reach/moves.py}; log \texttt{moves\_log.txt}), kinks with a stuck-edge background on their left are glued next to the edge of configurations that agree with the background to the right of the edge, and the glued configuration is followed in time; a move on the last $k$ cells and the exterior cell is recorded when the result is a configuration of the same kind changed only on these cells. The moves are tabulated on all $4^{k+1}$ states of these cells, with the cells further left drawn at random in $4$ independent contexts, a state being used only if all four give the same result, and their orbits are computed by union--find. The runs are: rule $87$ (class of $31$), background $(0,0)$, kinks of speed $1/3$, $k\le6$; rule $162$ (class of $13$), the background with rows $x=(00)^\infty$, $y=(01)^\infty$, speed $1/3$, $k\le4$; rule $25$, the background with rows $x=(01)^\infty$, $y=(10)^\infty$, speed $1$, $k\le4$; rule $74$ (class of $26$), background $(0,0)$, speed $1$, $k\le4$; rule $114$ (class of $58$), background $(1,1)$, speed $1$, $k\le4$; rule $191$, background $(0,0)$, speed $1/2$, $k\le4$; and rule $159$, background $(0,0)$, speed $1/2$, where no kink gives a move. As controls, the program reproduces the orbits of the moves described in \cite{PaperA} for $110$ ($k\le7$), $42$ and $46$ ($k\le8$).

\emph{Window graphs} (Remark~\ref{cl:rem:eigen}). The graphs $W^{(0,1)}_M$, $W^{(1,1)}_M$ and $W^{(-1,1)}_M$ are built from the cell rule for $M\le6$, for the smallest rule of each of the $88$ classes, and their strongly connected components are computed (\texttt{code/windows/window\_graphs.py}; output \texttt{window\_graphs\_M0-6.txt}).

\subsection{Confirmations of results proved by hand}\label{app:confirm}

The following computation checks a result that is proved by hand. It is not used in a proof.

\emph{Rule $110$R, no walls} (Proposition~\ref{cl:prop:nowalls}). For (1), over all $4^4$ values of $c_{[r-3,r]}$, the first $t$ at which the right context matters is $1$, $2$ and $3$ for $192$, $48$ and $16$ of them. For (2), over all windows and all three changes of a cell, $147{,}456$ pairs have a difference that stays in its column through $t=4$ and none through $t=5$, so $t\le5$ suffices (\texttt{code/searches/nowalls110.py}, exhaustive).

\subsection{Running times and reproducibility}\label{app:repro}

The censuses were run on one core of a machine with 32\,GB of memory under Linux, using Python~3.12 or later, NumPy, SciPy and PySAT~1.9 with CaDiCaL~1.5.3, MiniSat~2.2, Glucose~4 and MiniCard. Table~\ref{tab:times} lists instance sizes and times. The solution sets for rules $110$, $30$, $54$ and $90$ were each produced in two runs, with identical results. The radius-$4$ instances were solved with MiniSat~2.2, which needed about 5\,GB; CaDiCaL exceeded 12\,GB on the same instance. The 54R radius-$2$ invertible census used MiniCard's native cardinality constraints for the balance condition.

\begin{table}[h]
\caption{Instance sizes and running times (encode / enumerate / classify, in seconds).}
\label{tab:times}
\centering
\begin{tabular}{lrrrl}
\toprule
instance & variables & clauses & solutions & time (s)\\
\midrule
110R, 30R, $R=1$ & $128$ & $1.8\times10^4$ & 7 & $<1$\\
110R, 30R, $R=2$ & $2{,}048$ & $2.9\times10^5$ & 15 & $0.5$ / $0.3$ / $1$\\
110R, 30R, $R=3$ & $32{,}768$ & $4.7\times10^6$ & 27 & $7$ / $15$--$22$ / $33$\\
110R, 30R, $R=4$ (MiniSat) & $524{,}288$ & $7.5\times10^7$ & 43 & $130$--$160$ / $70$ / $47$\\
54R, $R=2$, balanced & $2{,}048+4{,}096$ & $2.9\times10^5$+card. & 23 & $3$ / $11$ / --\\
90R, $R=3$, elimination & $4{,}096$ classes & $2.6\times10^5$ eqs & $2^{65}$ & $4$\\
\bottomrule
\end{tabular}
\end{table}
